% Mục A.1. Phụ trách: TV1. Chuyển từ bản Word của TV1 (26/11/2026), đã rút gọn.
\phutrach{TV1 -- mục A.1}

\subsection{Đặc trưng dữ liệu vận đơn}
Bảng \code{shipments} của đề tài mô phỏng một hệ thống vận đơn với 100.000 cửa hàng (số đơn mỗi cửa hàng theo phân phối Zipf, tham số 1,1), 200 kho trung chuyển và 80\% đơn ở trạng thái đã giao. Trên cùng bảng này, hệ thống phải đáp ứng đồng thời ba nhóm yêu cầu:
\begin{itemize}
  \item \textbf{Giao dịch ACID} khi tạo đơn, cập nhật trạng thái và điều chỉnh số tiền thu hộ (\code{cod\_amount}).
  \item \textbf{Tra cứu qua nhiều chỉ mục phụ}: tra điểm theo \code{tracking\_code}, tra theo cửa hàng và thời gian tạo (\code{sender\_id}, \code{created\_at}), lọc theo khoảng \code{cod\_amount}.
  \item \textbf{Thống kê theo kho và trạng thái} (\code{dest\_hub\_id}, \code{status}), chạy song song với luồng cập nhật.
\end{itemize}

\subsection{Các nhóm hệ quản trị và hạn chế}
\textbf{RDBMS trên một máy.} MySQL (InnoDB) hay PostgreSQL hỗ trợ đầy đủ giao dịch ACID, SQL chuẩn và nhiều loại chỉ mục phụ, kể cả covering index trên B+Tree. Tuy nhiên toàn bộ bảng và chỉ mục nằm trên một node, nên dung lượng, bộ nhớ đệm và khả năng chịu lỗi bị giới hạn bởi phần cứng của máy đó. Khi tải ghi vượt một máy, phải chia dữ liệu thủ công (sharding) ở tầng ứng dụng hoặc tầng trung gian, khiến giao dịch và chỉ mục toàn cục phức tạp hơn.

\textbf{Wide-column NoSQL.} Cassandra hay HBase ghi nhanh và mở rộng ngang tốt. Đổi lại, nhóm này không có giao dịch ACID nhiều dòng, nhiều bảng theo mô hình quan hệ, chỉ mục phụ hạn chế và không hỗ trợ JOIN, nên khó phục vụ các truy vấn lọc theo nhiều thuộc tính độc lập của bài toán vận đơn.

\textbf{NewSQL phân tán.} TiDB giữ mô hình quan hệ, giao thức MySQL và giao dịch ACID ở tầng trên, còn tầng lưu trữ là các cặp Key-Value tự động chia thành Region trên nhiều node TiKV~\cite{pingcap_arch}. Bảng~\ref{tab:nhom-hqt} đối chiếu ba nhóm với yêu cầu của bài toán vận đơn. Cái giá của cách tổ chức này đối với chỉ mục phụ là câu hỏi mà đề tài đo.

\begin{table}[!htp]
\centering\bangfont
\caption{So sánh các nhóm hệ quản trị với bài toán vận đơn}\label{tab:nhom-hqt}
\begin{tabular}{|L{2.6cm}|L{3.9cm}|L{3.9cm}|L{3.9cm}|}
\hline
\dongdau \textbf{Tiêu chí} & \textbf{RDBMS một máy} & \textbf{Wide-column} & \textbf{NewSQL (TiDB)} \\ \hline
Giao dịch ACID & Đầy đủ, trên một máy & Hạn chế giao dịch nhiều dòng, nhiều bảng & ACID phân tán nhiều dòng (Snapshot Isolation) \\ \hline
Chỉ mục phụ & Đầy đủ trên B+Tree, có covering index & Hạn chế, khó truy vấn theo nhiều thuộc tính & Có, kể cả covering index, trên nền KV phân tán \\ \hline
Mở rộng ngang & Giới hạn bởi một máy; ghi nhiều phải sharding & Tốt cho cả lưu trữ và tải ghi & Mở rộng trực tuyến, tách tầng SQL và tầng lưu trữ \\ \hline
Chịu lỗi & Phụ thuộc một máy nếu không dựng thêm bản sao & Nhờ bản sao trên nhiều node & Tự bầu leader khi mất thiểu số bản sao (Multi-Raft) \\ \hline
\end{tabular}
\end{table}
